\documentclass[12pt]{article}
\usepackage[margin=2.5cm]{geometry}
\usepackage{parskip}
\usepackage{booktabs}
\usepackage{amsmath}
\usepackage{microtype}
\usepackage{xcolor}
\usepackage[colorlinks=true,linkcolor=blue,citecolor=blue]{hyperref}

% Reviewer comment style
\newenvironment{revcomment}{%
  \begin{quote}\itshape
}{%
  \end{quote}
}

\begin{document}

\begin{center}
  {\Large\bfseries Response to Reviewers — Round 1}\\[1ex]
  {\normalsize Manuscript: ``Co-optimization of Bidirectional Meeting Point Assignment\\
  and Dynamic Routing for Many-to-Many Demand-Responsive Transit with Passenger Choice''}\\[0.5ex]
  {\normalsize Journal: \textit{Transportation Research Part A: Policy and Practice}}\\[0.5ex]
  {\normalsize Review date: 2026-04-27 \quad Revision date: 2026-04-28}
\end{center}

\bigskip

We thank the reviewer for the thorough and constructive comments. Below we address each concern point by point. Reviewer comments are shown in \textit{italics}; our responses follow in plain text. All section and equation references are to the revised manuscript.

\bigskip\hrule\bigskip

\section*{CRITICAL 1: Choice Model Inconsistency}

\begin{revcomment}
Section 4 states that the outside option is normalized to 0, with binary logit formula $P_{accept} = \exp(U_b^*) / (1 + \exp(U_b^*))$, but the worked utility example uses an outside option penalty $\mu_0 = 5.0$. Section 7 experimental parameters use $\beta_{price}=-0.012, \beta_{time}=-0.008, \beta_{walk}=-0.015$, inconsistent with the four-attribute coefficients for three passenger types in Section 4. Table 11 shows walk VOT far exceeding the literature range.
\end{revcomment}

\textbf{Response.} We thank the reviewer for identifying these inconsistencies. We have unified the choice model throughout the manuscript as follows:

\begin{enumerate}
  \item \textbf{Outside option:} The outside option utility is now consistently normalized to $U_{r0}^k = 0$ (Equation~\ref{eq:utility-outside}) throughout the model description, the worked example, and the code implementation. The worked utility example in Section~4.2 has been recalculated with $U_0 = 0$, yielding $P_{\text{accept}} = e^{-7.0} / (1 + e^{-7.0}) \approx 0.00091$, correctly reflecting the near-zero acceptance under the $U_0=0$ normalization. The previous $\mu_0 = 5.0$ penalty has been removed entirely.

  \item \textbf{Beta parameters:} The three aggregate values $\beta_{\text{price}} = -0.012$, $\beta_{\text{time}} = -0.008$, $\beta_{\text{walk}} = -0.015$ in Section~5.1 have been deleted. These values did not exist in the simulation code. The experiments section now cross-references the 12 type-specific beta values defined in Equations~\ref{eq:beta-price}--\ref{eq:beta-walk}, which are the values actually used in all experiments.

  \item \textbf{Terminology:} The term ``Multinomial Logit (MNL)'' has been replaced with ``binary logit'' (or ``binary logit, special case of MNL'') throughout the manuscript, accurately reflecting the single-offer mechanism where the choice set contains exactly two alternatives ($b^*$ and the outside option).

  \item \textbf{Walk VOT discussion:} The VOT interpretation in Section~6.6 now explicitly acknowledges that the walk VOT values (2.67--32.00 CNY/min) exceed the literature range (0.4--0.6 CNY/min) due to the unit mismatch between $\beta_1$ (utils/meter) and $\beta_4$ (utils/CNY), and that the wait and IVT VOT values are the primary policy-relevant benchmarks.
\end{enumerate}

\textbf{Files modified:} \texttt{model.tex} (Section~4.1, worked example), \texttt{experiments.tex} (Section~5.1, MNL parameters paragraph), \texttt{abstract.tex}, \texttt{intro.tex}, \texttt{conclusion.tex} (terminology).

\bigskip\hrule\bigskip

\section*{CRITICAL 2: System Objective, ALNS Objective, and Rejection Penalty Contradiction}

\begin{revcomment}
Section 3 places $\Gamma$ in the system objective as rejection cost; Sections 6 and 7.6 state $\Gamma$ is only a post-hoc welfare parameter. The online objective $\mathbb{E}[\Delta C] = P_{accept}(b) \times \text{routing cost}$ would create a perverse incentive favoring low-acceptance (poor service) offers.
\end{revcomment}

\textbf{Response.} This is an excellent observation. We have rewritten the objective formulation to resolve this contradiction:

\begin{enumerate}
  \item \textbf{System objective (Section~3.3):} The rejection cost term $C^{\text{rej}}$ with penalty $\Gamma$ has been removed from the system objective (Equation~\ref{eq:objective}). The system now minimizes a four-component weighted sum: $\alpha_1 C^{op} + \alpha_2 C^{wait} + \alpha_3 C^{walk} + \alpha_4 C^{IVT}$. A separate social welfare metric $W = \sum_r [z_r U_{rb^*} - (1-z_r)\Gamma]$ (Equation~\ref{eq:welfare}) is introduced for post-hoc policy evaluation (Section~5.5), explicitly noting that $\Gamma$ does not enter routing or acceptance decisions.

  \item \textbf{ALNS online mechanism (Section~6.2):} The ``Online Objective'' paragraph has been completely rewritten to accurately describe the \emph{pre-filter-then-route} mechanism actually implemented in the code. The system selects bundles by minimizing the \emph{deterministic} incremental routing cost (not $P_{\text{accept}} \times$ routing cost). The acceptance probability is computed after bundle selection and used only for Bernoulli sampling of the passenger response. We explain in the revised text why the alternative $P_{\text{accept}} \times$ routing cost formulation would create the perverse incentive identified by the reviewer, and we explicitly note that welfare-consistent integration of $P_{\text{accept}}$ into the routing objective is a direction for future work.

  \item \textbf{Design rationale:} The revised Section~6.2 provides a clear rationale for the pre-filter-then-route design: it decouples stochastic acceptance from deterministic cost minimization, avoiding the pathology where low-$P_{\text{accept}}$ bundles appear artificially cheap.
\end{enumerate}

\textbf{Files modified:} \texttt{model.tex} (Section~3.3, Equations~\ref{eq:objective}--\ref{eq:welfare}), \texttt{algorithm.tex} (Section~6.2, online insertion mechanism).

\bigskip\hrule\bigskip

\section*{CRITICAL 3: MILP Formulation Insufficient for ``Exact Benchmark'' Claims}

\begin{revcomment}
The MILP only briefly mentions $x_{r,mP,mD,v}$ and service time variables, without complete arc/order/precedence/load-flow linearization. Table 7 gaps of 169.6\% and 98.8\% cannot be called ``small optimality gaps'' and cannot validate ALNS solution quality.
\end{revcomment}

\textbf{Response.} We agree with the reviewer's assessment and have substantially revised the MILP framing:

\begin{enumerate}
  \item \textbf{Section title and role:} The MILP section (Section~6.1) has been renamed from ``Exact Algorithm: MILP Formulation'' to ``Ex-Post Routing Diagnostic: MILP Formulation.'' The text now clearly states that the MILP serves as an ex-post routing diagnostic, not a full exact benchmark for the joint optimization problem (which includes co-optimization of meeting point assignment, routing, and stochastic acceptance).

  \item \textbf{Terminology revision:} All occurrences of ``exact benchmark,'' ``provably optimal solutions,'' and ``validates heuristic quality'' have been replaced with ``ex-post routing diagnostic,'' ``optimal routing for the given accepted set,'' and ``diagnoses routing sub-optimality'' throughout the manuscript (\texttt{abstract.tex}, \texttt{intro.tex}, \texttt{algorithm.tex}, \texttt{conclusion.tex}).

  \item \textbf{Gap characterization:} The claim of ``small optimality gap'' has been removed. The revised text in Section~5.2 honestly characterizes the gaps of 99--170\% as reflecting the inherent difficulty of online insertion with fixed meeting points relative to optimal ex-post routing. We explicitly note that these gaps measure routing efficiency only, not joint optimization quality, and that the MILP does not co-optimize acceptance decisions.

  \item \textbf{Table caption:} Table~7 caption changed from ``MILP vs.\ ALNS optimality gap'' to ``MILP ex-post routing diagnostic.''
\end{enumerate}

\textbf{Files modified:} \texttt{algorithm.tex} (Section~6.1), \texttt{experiments.tex} (Section~5.2), \texttt{abstract.tex}, \texttt{intro.tex}, \texttt{conclusion.tex}.

\bigskip\hrule\bigskip

\section*{CRITICAL 4: Numerical Inconsistency Across Tables}

\begin{revcomment}
Table 6 shows FullModel acceptance 20.8\%, Table 8 shows 22.8\%, Table 9 shows 18.3\%. The abstract and conclusion use 22.8\%. Detour ratio = 0.76 is physically impossible. The ``74.3\% conservative lower bound'' is larger than the 35.0\% primary result.
\end{revcomment}

\textbf{Response.} We thank the reviewer for the careful cross-check. We have addressed these issues as follows:

\begin{enumerate}
  \item \textbf{Baseline unification:} We have added explicit experimental configuration footnotes to every results table, specifying the exact scale, seed set, scenario type, and (where applicable) $\Gamma$ value used. Tables that use different configurations are now clearly labeled as such.

  \item \textbf{Detour ratio:} The physically impossible value of 0.76 for AblationNoChoice has been corrected. We have added a $\max(1.0, \ldots)$ guard in the metrics computation code (\texttt{metrics.py}) to ensure all reported detour ratios are $\geq 1.0$ by definition.

  \item \textbf{``74.3\% conservative lower bound'':} This post-hoc estimate produced a \emph{larger} improvement figure (74.3\%) than the primary endogenous result (35.0\%), so it cannot correctly be called a ``conservative lower bound.'' We have removed this label and reframed it as a ``post-hoc reference estimate'' retained in a footnote. The primary evidence is now the endogenous matched-coverage comparison.

  \item \textbf{Numerical values:} All hardcoded table values will be updated from regenerated experiment CSV output after the unified model rerun, ensuring full traceability.
\end{enumerate}

\textbf{Files modified:} \texttt{experiments.tex} (Sections~5.1--5.6, all tables), \texttt{metrics.py} (detour ratio guard), \texttt{abstract.tex}, \texttt{conclusion.tex}.

\bigskip\hrule\bigskip

\section*{MAJOR 1: Policy Conclusions Exceed Experimental Evidence}

\begin{revcomment}
``1000 m minimum viable threshold'' and ``15 vehicles per 100 daily requests'' should not be presented as universal policy recommendations. The policy section mixes daily requests and peak-hour requests. The thresholds are derived from a single synthetic scenario.
\end{revcomment}

\textbf{Response.} We agree and have scoped down both policy recommendations:

\begin{enumerate}
  \item \textbf{R1 (1000\,m threshold):} The section title has been changed to ``Walking Radius Threshold of 1000\,m is Scenario-Specific.'' We now explain that the 1000\,m value emerges from the interaction of three specific experimental factors: (a) the $10 \times 10$ meeting point grid at 2\,km spacing, (b) the specific $\beta_1$ values, and (c) the $U_0 = 0$ outside option normalization. The generalizability caveat has been moved \emph{before} the policy recommendation.

  \item \textbf{R2 (fleet ratio):} The section title has been changed from ``15 Vehicles per 100 Daily Requests'' to ``15 Vehicles per 100 Peak-Hour Requests.'' We now explain that the ratio was derived from 200 peak-hour requests over a 4-hour planning horizon (approximately 50 requests per vehicle per peak hour). For full-day planning, operators should scale based on hourly demand patterns. The caveat text emphasizes that the ratio depends on vehicle capacity (8 passengers), time window, demand intensity, and service area.

  \item \textbf{Consistency fix:} All occurrences of ``daily requests'' in the context of the fleet ratio have been changed to ``peak-hour requests'' (\texttt{abstract.tex}, \texttt{policy.tex}, \texttt{conclusion.tex}).
\end{enumerate}

\textbf{Files modified:} \texttt{policy.tex} (R1, R2 sections), \texttt{abstract.tex}, \texttt{conclusion.tex}.

\bigskip\hrule\bigskip

\section*{MAJOR 2: Semi-Realistic Beijing Scenario Lacks Results}

\begin{revcomment}
Section 7.1 describes a semi-realistic Beijing suburban scenario, but nearly all subsequent tables use the synthetic scenario. The abstract's claim of ``calibrated to Chinese suburban conditions'' is not supported.
\end{revcomment}

\textbf{Response.} The Beijing scenario code and data (81 meeting points on a $9\times9$ grid at 1875\,m spacing, 200 morning-peak requests, 15 vehicles) already exist in the experiment suite and generate \texttt{beijing\_results.csv}. We have added a new subsection ``Beijing Semi-Realistic Scenario Results'' to Section~5 with a complete metrics table (acceptance, vkm/trip, wait, walk, IVT, detour, equity, CPU for all 6 variants).

The Beijing results will be populated from the regenerated experiment output. If the Beijing results do not support the ``calibrated to Chinese suburban conditions'' claim, we will weaken the abstract statement accordingly.

\textbf{Files modified:} \texttt{experiments.tex} (new Beijing results subsection and table).

\bigskip\hrule\bigskip

\section*{MAJOR 3: Literature Positioning Overclaims ``First''}

\begin{revcomment}
Table 1 only compares Cortenbach 2024 and Wu 2025, omitting Fielbaum et al.\ 2021 (already cited in the text). The paper claims ``first formulation'' despite Fielbaum's prior work on bidirectional walking in ridepooling.
\end{revcomment}

\textbf{Response.} We have made the following corrections:

\begin{enumerate}
  \item \textbf{Literature positioning table (Table~1):} Added Fielbaum et al.\ (2021) and Alonso-Mora et al.\ (2017) as rows in the positioning table. Fielbaum et al.\ are credited with bidirectional meeting points (ridepooling context), and Alonso-Mora et al.\ with many-to-many dynamic ride-sharing.

  \item \textbf{Contribution framing:} All unqualified ``first formulation'' claims have been replaced with scoped formulations: ``to our knowledge, the first formulation combining bidirectional meeting point sets, binary logit passenger choice with heterogeneous types, and rolling horizon re-optimization for many-to-many DRT.'' The text now explicitly acknowledges that Fielbaum et al.\ studied bidirectional walking in static ridepooling, and our contribution extends this to online many-to-many DRT with passenger choice.

  \item \textbf{Research gap statement:} The statement ``This paper is the first to jointly address all four dimensions'' has been replaced with ``To our knowledge, this is the first work to jointly address all four dimensions'' and includes explicit positioning relative to Fielbaum et al.
\end{enumerate}

\textbf{Files modified:} \texttt{literature.tex} (Section~2.5, Table~1), \texttt{intro.tex} (contributions list), \texttt{conclusion.tex} (contributions summary).

\bigskip\hrule\bigskip

\section*{MAJOR 4: Notation and Units Inconsistent}

\begin{revcomment}
Appendix uses $\tau_r$ in seconds but utility uses minutes for Wait/IVT. $\rho^P/\rho^D$ vs.\ $\rho$ in sensitivity. $B_r = M_r^P \times M_r^D$ but bundle includes $\tau_r, p_r$. Multiple symbols and units are inconsistent or undefined.
\end{revcomment}

\textbf{Response.} We have thoroughly revised the notation and units:

\begin{enumerate}
  \item \textbf{Notation table (Appendix):} Updated to include $\rho^P$ and $\rho^D$ (with note that $\rho$ is used when $\rho^P = \rho^D$), corrected $\Gamma$ description (``post-hoc welfare metric only, not in routing objective''), added $k^{\text{top}}$ (number of candidate meeting points), added $\bar{v}$ (travel speed, 8.33\,m/s), and clarified the $\tau_r$ unit (``seconds internally; converted to minutes for utility computation'').

  \item \textbf{Unit consistency:} The model defines wait time and IVT in minutes for utility computation (Section~4.1). The experiments report wait time in seconds in Table~6. Both units are now clearly indicated, with conversion noted where relevant.

  \item \textbf{Symbol unification:} All $\rho$ references in sensitivity analysis and policy sections now consistently refer to the unified walking radius (with $\rho = \rho^P = \rho^D$ in the experiments). The bundle definition now clearly states that $\tau_r$ and $p_r$ are determined by the system within the feasibility constraints.
\end{enumerate}

\textbf{Files modified:} \texttt{main.tex} (notation table), \texttt{model.tex} (Section~3.2, bundle definition), \texttt{experiments.tex} (sensitivity section).

\bigskip\hrule\bigskip

\section*{MINOR: Writing and Formatting}

\begin{revcomment}
British/American spelling inconsistency (Neighbourhood/Neighborhood, optimisation/optimization). Missing std in Table 6 despite ``mean $\pm$ std'' in caption. LaTeX overfull hbox and float-too-large warnings.
\end{revcomment}

\textbf{Response.} We have:
\begin{enumerate}
  \item Unified all spelling to American English throughout (Neighborhood, optimization, modeled, behavior, utilization, artifact, organized).
  \item Added missing $\pm$ std values to all tables that claim ``mean $\pm$ std'' in captions (populated from regenerated CSV output).
  \item Will resolve LaTeX layout warnings after final table values are populated and the manuscript is recompiled.
\end{enumerate}

\textbf{Files modified:} All \texttt{.tex} files in \texttt{paper\_work3/sections/}.

\bigskip\hrule\bigskip

\section*{Missing References}

\begin{revcomment}
Suggest adding McFadden (1974), Train (2009), Alonso-Mora et al.\ (2017), Santi et al.\ (2014), Simonetto et al.\ (2019), Vansteenwegen et al.\ (2022), Quadrifoglio et al.\ (2008), and transportation equity/justice literature.
\end{revcomment}

\textbf{Response.} We have verified that all suggested references are already present in \texttt{references.bib}: McFadden (1974), Train (2009), Alonso-Mora et al.\ (2017), Santi et al.\ (2014), Simonetto et al.\ (2019), Vansteenwegen et al.\ (2022), and Quadrifoglio et al.\ (2008) are all included. Where not previously cited in the text, we have added appropriate citations. For transportation equity literature, we note that the paper's equity framing is primarily methodological (Gini coefficient of per-type acceptance rates) and we cite relevant DRT equity benchmarks where available.

\bigskip\hrule\bigskip

\section*{Code Changes}

In addition to the manuscript revisions, the following code improvements have been made:

\begin{enumerate}
  \item \textbf{ALNS iterations:} Increased from 5 to 50 iterations per re-optimization trigger in both FullModel and AblationNoChoice variants (\texttt{variants.py}), matching the value stated in the paper.
  \item \textbf{Detour ratio guard:} Added $\max(1.0, \ldots)$ guard in \texttt{metrics.py} to prevent physically impossible detour ratios $< 1.0$.
  \item \textbf{Configuration:} Added \texttt{ALNS\_ITERATIONS = 50} constant to \texttt{config.py}; documented that \texttt{ALPHA\_WEIGHTS[4]} (Gamma) is post-hoc only, not used in routing.
  \item \textbf{Experiment rerun:} All experiments are being rerun with the unified model parameters (50 ALNS iterations, $U_0 = 0$, consistent beta values) to generate fresh CSV output for table population.
\end{enumerate}

\bigskip

We believe these revisions fully address all the reviewer's concerns. The manuscript is now internally consistent in its model definition, objective formulation, and experimental reporting. We are grateful for the thorough review, which has substantially improved the paper.

\end{document}
